\subsection{Whole \(D_{17}\): fresh geometry before any field census}
\label{app:arithD17geometry}

Use ACTUAL unit \(D_{17}\) roots, long spine one
through fifteen, forks sixteen/seventeen, and
\[
 x_1>\cdots>x_{16}>|x_{17}|,\qquad\|x\|<1 .
\]
The signed coordinate is retained.
The legal generic \(M=11\) window gives strong
six through twelve.
At thirteen/fourteen the COMPLETE right D4/D3
lists leave only the pure own-five/fork-three
or own-five/end-three packets.
Here completeness uses the SAME actual finite
role list printed in Section~\ref{app:arithD13rankthree},
with the FRESH indices thirteen/fourteen.
The D4 long-three and D3 middle-three roles have
residual \(1/2\), shift \(1/2\);
fork/end-three has residual \(1/2\) or \(5/8\),
shift \(1/4\).
D3 end-four has residual \(1/4\), shift
\(1/(2\sqrt2)\); sealed end-five has residual
\((5-3\varphi)/8\), shift \(\varphi/4\).
The \(D_{17}\) bridge residual is \(1/(2i)\);
the \(A_{17}\) end-three bridge at fourteen
has residual \(1/8+1/28=9/56\).
Their positive-own numerators make them strong,
and their own-zero cycles are affine.
For single-fork/end-three, own zero, three,
four and six retain the actual affine \(E_8\),
cycle, \(F_4\) and \(6,3\) principals.
The own-three cycle has FIVE actual nodes at
D4 and FOUR at D3 end-three; it is not
replaced by a triangle.
The far-helper arithmetic lemma removes the
other nonisolated possibility.
The original LEFT isolated-core/H4/A2/singleton
proof applies with \(2i\ge26\), so its previously
proved \(2i\ge18\) loss margins suffice.
The remaining fork-helper raw Schur constant
is ACTUAL seventeen; the complete norm/trace
argument leaves only its golden own.
There is no spectator or extra long root here.
Their actual LEFT \(H_4\) points have divided
displacement \(V\), flat one through \(K=i-3\),
\(V^2<i\), and pressure
\[
 x\cdot V>G>\sqrt{i-\kappa/(\varphi+5/8)}-.091,\quad
 \kappa=1/2,\ 5/8 .
\]
Every point uses positive actual walls, positive
future time and all original-wall signs.
The relevant rows are
\[
\begin{array}{c|c|c|c|c}
 i&K&\mu&G>&f_K<\\ \hline
13&10&.16&3.48&2.732\\
14&11&.15&3.61&2.886 .
\end{array}
\]
The radical seals are
\(13-25/112>3.571^2\),
\(14-125/448>3.701^2\).
The INDIVIDUAL terminal prices exceed
\(.16,.15\), in both initial-five references.
The residual \(V-P_K\) vanishes on the whole
reference support and has norm \(<3\).

If five is weak, its actual endpoint reference
has energies \(L_{10}>.8870,L_{11}>.9100\).
The two paid squares are
\[
 .8870+.32(3.48-2.732)-.0256(3)>1.0495,\qquad
 .9100+.30(3.61-2.886)-.0225(3)>1.0597 .
\]
If five is strong, first classify the ACTUAL
head: the retained \(D_{16}\) admits only isolation
or simple three at its long endpoint two.
The nonisolated row gives \(x_1>.643\), hence
legal first-five gain \(>.0994\).
An isolated head has small owns/nulls excluded
as before, including five by
Section~\ref{app:arithgoldenlongstar}.
For finite own at least seven, its actual finite
head point gives gain \(>.0135\) at cost \(.1\)
from \(P_5\).
Explicitly for \(1<u<1.232\),
\[
 \eta=H-.361-(R-.2)/4>.0135 .
\]
On \(u\in[1,1.1],[1.1,1.2],[1.2,1.232]\),
the respective rational seals are
\[
 H>.6708-.1871,\quad .63245-.204125,\quad .6196-.2,
 \qquad R\le.5,.415,.36 .
\]
The middle second-radical square has maximum
\(1/24\); all bounds are positive squarings.
They give gains \(>.0477,.013575,.0186\).

For a dotted head use its ACTUAL positive core,
\(Z=(1-u/2)e_1+(u/2)e_2\), \(0<u<1\).
Paying \(.36\) at \(P_5\), cost \(.072\),
gives gain
\[
                .72(H-b)-.1296(R-1/5)>.04 .
\tag{D17H}\label{eq:arithD17dotgain}
\]
Here is the full continuous seal rather than a
borrowed actual-\(I\) theorem.
On \([0,1/2],[1/2,3/4],[3/4,1]\), use lower
chords for \(\sqrt{2-u}\)
\[
1.4142-.379u,\quad1.4381-.4268u,\quad1.4721-.4721u
\]
and upper chords for \(R\)
\[
1-.75u,\quad.8125-.375u,\quad.625-.125u.
\]
With \(.5091<9\sqrt2/25<.5092\), the three
gain bounds reduce to
\[
\begin{array}{l}
 .35636922-.0958u-.5092\sqrt{u(1-u)},\\
 .39283671-.16873u-.5092\sqrt{u(1-u)},\\
 .43444611-.2242u-.5092\sqrt{u(1-u)} .
\end{array}
\]
The circle inequality and squared-root upper
bounds \(.5182,.5365,.5564\) give minima
\(>.04936922,.04022171,.04414611\).
No \(I\) wall or energy enters.

Every head residual is zero-sum on the first
five and orthogonal to the H4 residual beyond
\(K=10/11\).
The bare strong-five H4 bounds are
\(.98692,.99786\).
Adding even the smallest LEGALLY funded head
gain \(>.0135\) excludes both.
All six through fourteen are consequently strong
BEFORE deciding five or the final head.

Weak five with finite own gives actual
\(x_1>.65\); its rebuilt six-through-fourteen
reference exceeds one.
Its dotted case is precisely the complete
minimal-degree norm/trace argument already
proved for the SAME inverse entries in
Section~\ref{app:arithB17physical}.
Thus five is strong.
A nonisolated head now gives
\(E_{14}+.0994>1\), so it is isolated.
The legal graft at \(x_1\ge.632\) gives
\(E_{14}+(5/4)(.632-.361)^2>1.00020125\).
Therefore \(x_1<.632\).
The original head raw normalization gives
\[
 r_H^2=2U,\quad\text{spatial }r_H=Ue_1,\quad
 H_0^2=U(U-2),\quad K=\mathbb Q(U),\quad
               2<U<11,\quad0<U_\sigma<2 .
\]
Finite heads remain possible at this stage.

For completeness the conjugate interval itself
bounds their COMPLETE native alphabet.
For \(U=2+2\cos(2\pi/m)\), a coprime exponent
\(1<a<m/4\) gives a forbidden conjugate \(>2\).
For odd \(m\ge9\) use two.
For even \(m\ge14\) not divisible by three use three.
For \(m=6k\), \(k=4,5\), use five/seven;
for larger even \(k\) use \(k-1\), or \(k-3\)
if \(k\equiv1\pmod3\);
for larger odd \(k\) use \(k-2\), or \(k-4\)
if \(k\equiv2\pmod3\).
Each is coprime and below \(3k/2\).
The remaining labels are
\(3,4,5,6,7,8,10,12,18\).
The prior null/golden exclusions leave exactly
seven/eight/ten/twelve/eighteen, with the
parameters already printed for \(B_{17}\).

For either actual signed fork direction
\(V_f=P_{16}\pm e_{17}\), if \(C_f\ge13\),
the ORIGINAL edge gives
\(x\cdot V_f\ge\sqrt{102/7}>3.817\).
Use the SAME head gain and a separately funded
\(.13(V_f-P_{14})\), whose residual norm is three
and cross with the head/reference is zero.
Then
\[
 \|x\|^2>.9084+.0135+.26(3.817-3.311)-.0169(3)
                                  =1.00276>1 .
\]
Thus BOTH actual fork coshes are \(<13\),
without equating their parameters.

\subsection{The complete actual \(D_{17}\) fork table}
\label{app:arithD17forktable}

Fix either genuine fork \(f\).
Deleting it retains ACTUAL \(A_{16}\).
Finite classification gives isolation, helper three
at \(j=1,2,15,16\), or helper four at \(j=1,16\).
For three put \(t=2\gamma\), raw cap square two,
white product \(-1\), own product \(-t\).
For four put \(a=2\sqrt2\gamma\), raw square one,
white product \(-1\), own product \(-a/2\).
Actual cycles make the parameter integral,
and for four also \(a^2/2\) integral.
The original Lorentzian principal gives ACTUAL
field equality with that parameter.
The COMPLETE raw time and actual edge table is
\[
\begin{array}{c|c|c}
 \text{helper}&H_c^2&C_f\\ \hline
3@1&(17t^2+4t-4)/4&(17t+2)/(6\sqrt2)\\
3@2&(17t^2+8t)/4&(17t+4)/4\\
3@15&(17t^2+60t+52)/4&(17t+30)/4\\
3@16&(17t^2+30t+9)/4&(17t+15)/(6\sqrt2)\\
4@1&(17a^2+8a)/16&(17a+4)/4\\
4@16&(17a^2+60a+52)/16&(17a+30)/4 .
\end{array}
\]
The near endpoint THREE formula differs from
the B17 short formula and is not substituted
for it.
Every other conjugate has negative displayed
time square.
Every null below uses UNIT ACTUAL cap normals;
its positive time and all original-wall signs
follow termwise.

The \(3@15\) own-zero kernel is
\(n_{14}+2n_{15}+n_{16}+n_{17}+c\);
own three has its triangle.
All other owns have \(t\ge\sqrt2\), \(C_f>13\).
The \(4@16\) own-zero kernel is
\(n_{14}+2n_{15}+2n_{16}+n_{17}+\sqrt2c\);
every positive \(a\ge\sqrt2\) gives \(C_f>13\).
At \(3@2\), own zero has zero time; positive
owns satisfy
\(1\le t<48/17\), \(-8/17<t_\sigma<0\).
Degree at least three has nonzero integer norm
magnitude \(<1\).
Degree two leaves norm \(-1\), trace one or two;
the golden root violates the conjugate strip,
leaving \(1+\sqrt2\).
Its raw time norm is \(-47/16\).
For a finite head the only compatible head is
eight, of time norm \(-4\), and the ratio
\(47/64\) is nonsquare.
For a DOTTED head, the complete actual
\(\sqrt2\) head strip gives time norms
\(-7,-7,-32,-63,-36\).
Prime forty-seven again has odd ratio valuation.
Thus this row is excluded under EVERY head.
At \(4@1\) the same norm strips leave only
\(a=1+\sqrt2\), but \(a^2/2=3/2+\sqrt2\)
is not integral.
Its own zero has zero time.
Every four non-end role is exhausted, with
finite, parallel and dotted own retained;
parallel everywhere has \(c+n_f\).
Far-end own zero has negative time, while
near-end own zero gives the forbidden rational
actual field.

For the two remaining end-three roles put
\(v=t\) far, \(v=t+1\) near.
Their raw times are
\[
 H_c^2=(17v^2\pm4v-4)/4,
\]
with upper sign far.
The COMPLETE strips are
\[
\begin{array}{c|c|c}
 \text{role}&\text{physical }v&\text{every other }v_\sigma\\ \hline
 \text{far}&1\le v<6.372&-.617<v_\sigma<.382\\
 \text{near}&2\le v<6.607&-.382<v_\sigma<.617 .
\end{array}
\]
The exact nonphysical endpoints are
\((\mp2\pm6\sqrt2)/17\), and the physical
upper endpoints are
\((78\sqrt2\mp2)/17\).
These exact signs are retained whenever a
safe decimal overstrip admits an extra row.

Let \(g_0\) be the ACTUAL raw head/helper
product.
Its actual triangle or complete old-chain
cycle makes it integral.
The common original time-line identities are
\[
 \begin{split}
 4g_0^2-4g_0U(t+2)+U^2(t+2)^2
                  &=U(U-2)(17t^2+4t-4),\\
 4g_0^2-4g_0U(t+1)+U^2(t+1)^2
                  &=U(U-2)(17t^2+30t+9).
 \end{split}
\]
Reducing modulo \(4\mathcal O_K\) gives the
necessary parity statement
\[
                              Uv^2\in2\mathcal O_K .
\tag{D17P}\label{eq:arithD17helperparity}
\]
It does NOT justify \(Uv/2\) integral at
ramified dyadic primes.
The physical and conjugate bounds belong to
actual \(K=\mathbb Q(v)\), not merely an
unexamined smaller parameter field.

\subsection{Whole \(D_{17}\): every finite head is excluded}
\label{app:arithD17finiteheads}

Assume one of the FIVE finite heads retained
in Section~\ref{app:arithD17geometry}.
Its actual field is the corresponding quadratic
or cubic field, so the complete fork table applies
with that field equality.

For heads seven/ten/eighteen, \(U\) is a dyadic
unit and two is unramified, by discriminants
forty-nine/five/eighty-one.
Parity~\eqref{eq:arithD17helperparity} makes \(v/2\)
integral.
In the far role its nonzero norm has magnitude
\(<3.186(.3085)^{d-1}<1\) for \(d=2,3\).
For the near role the cubic norm is
\(<3.304(.309)^2<1\).
In the remaining quadratic golden near case,
\(v/2=a+b\varphi\) has physical value \(<3.304\),
other value in \((-.191,.309)\).
The gap gives \(b\le1\), \(b>0\).
For \(b=1\) the interval requires
\(.427<a<.927\), no integer.
Both end-helper roles are impossible.

For head eight use \(K=\mathbb Q(s)\), \(s=\sqrt2\),
\(U=2+s\).
At its dyadic prime \(\operatorname{ord}(U)=1\),
\(\operatorname{ord}(2)=2\).
Parity requires \(s\mid v\).
The COMPLETE far strip has
\(t=1+s,3+2s\), neither divisible by \(s\).
The complete near strip has
\(v=2+s,3+2s\); only \(2+s\) survives,
giving \(t=1+s\).
Its raw time square has norm \(-23/4\),
with nonsquare ratio \(23/16\) to the actual
head time norm \(-4\).
For head twelve, \(U=2+\sqrt3\) is a dyadic unit
although two ramifies.
Parity requires \(v\) in its dyadic prime,
where \(\sqrt3=1\) in the residue field.
The COMPLETE far/near strips leave only
\(t=2+\sqrt3\), respectively \(v=2+\sqrt3\);
both have nonzero residue one.
Thus BOTH actual forks are isolated under
EVERY finite head.
The assertion is presently conditional on a
finite head, not carried into the dotted case.

For an isolated positive fork put \(V=4\gamma_f^2\).
Its ACTUAL raw root has square \(2V\),
sole old product \(-V\), time square
\[
              H_f^2=V(17V-8)/4,\quad
              K=\mathbb Q(V),\quad0<V_\sigma<8/17 .
\]
Own zero has no spatial wall part.
Own three has the actual affine \(E_8\) arms
one/two/five; own four has positive affine
vector
\(2n_{15}+2n_f+n_{14}+n_k+\sqrt2c\),
where \(k\) is the other fork.
Own six has its \(6,3\) path; parallel has
the actual null pair.
All are future feasible against EVERY original
wall.
For finite labels at least seven,
degree at least three contradicts
\(\mathrm N(V)<4(1/2)^2\le1\).
In degree two the strict strip forces norm one,
trace four, namely native twelve.
Native five is the other finite candidate.
The cubic/silver head fields cannot equal those
of these candidates.
At golden head ten the fork-five time norm
is \(-55/16\), giving nonsquare ratio \(11/16\)
to head norm \(-5\).
At head twelve the fork-twelve time norm
is \(-191/16\), ratio \(191/48\) to head norm
\(-3\), also nonsquare.
Both forks are DOTTED under every finite head.

The actual head/fork LONG cycle gives
\(g=(r_H,r_f)\) integral; the actual two-cycle
gives \(L=g^2/(UV)\) integral.
Their spatial product is \(UV/2\).
The SAME original time line gives
\[
 (g-UV/2)^2=U(U-2)V(17V-8)/4,\qquad
 L-g=4UV-17V/2-2U+4 .
\]
Hence \(W=V/2\) is integral, by subtracting
eight copies of \(V\) from the integral \(17V/2\).
This uses BOTH cycles without a dyadic-unit
assumption.
The actual fork bound gives
\[
                  2<W<676/17,\quad0<W_\sigma<4/17,\quad
                  H_f^2=W(17W-4).
\tag{D17W}\label{eq:arithD17forkstrip}
\]

For the three quadratic finite head fields,
the COMPLETE full-order strips and time-norm
ratios are
\[
\begin{array}{c|l|l}
 \text{head}&W&\mathrm N(H_f^2/H_0^2)\\ \hline
8&
 3+2\sqrt2,\ 10+7\sqrt2,\ 17+12\sqrt2&
 103/4,\ 383,\ 2007/4\\
10&
 2+3\varphi,\ 5+8\varphi,\ 7+11\varphi,\ 10+16\varphi&
 171/5,\ 919/5,\ 239,\ 5104/5\\
12&
 7+4\sqrt3,\ 14+8\sqrt3&
 647/3,\ 976 .
\end{array}
\]
For coverage, the positive coefficient bounds
are fourteen/seventeen/eleven, and for each
coefficient the small-conjugate width \(4/17<1\)
permits at most one integer constant term.
The physical bound removes the otherwise
small-conjugate \(20+14\sqrt2\).
For \(n=\mathrm N(W),T=\operatorname{Tr}(W)\),
\(\mathrm N(17W-4)=289n-68T+16\)
gives every printed ratio.
Each is a nonsquare by an odd denominator
three/five valuation or the nonsquare integer
numerator; \(976=16(61)\).
The original common \(K\)-time line therefore
excludes every row.

For finite heads seven and eighteen the FULL
rings are \(\mathbb Z[U]\).
Their polynomials have discriminants \(49,81\);
the former translated by four is seven-Eisenstein,
the latter is three-Eisenstein.
The distinct-valuation uniformizer argument
already proved for the cubic bases in
Section~\ref{app:arithA16anchors} excludes the
only possible index prime.
Both are irreducible modulo two, so two is inert.
The positive norm of \(W\) is \(<10816/4913<3\).
An even norm would be divisible by eight;
hence \(\mathrm N(W)=1\), \(W>289/16\).

Let \(\alpha,\beta,\chi\) be physical and two
small conjugates of the actual finite-head generator.
For head seven the exact root and derivative
bounds are
\[
 \begin{gathered}
3.246<\alpha<3.248,\quad1.554<\beta<1.556,\quad
 .198<\chi<.199,\\
5.149<f'(\alpha)<5.167,\quad|f'(\beta)|>2.289,\quad
 f'(\chi)>4.128,\\
1.752<\beta+\chi<1.754,\qquad\beta-\chi>1.355 .
 \end{gathered}
\]
They follow by polynomial signs and products
of the disjoint root gaps.
Writing \(W=A+BU+CU^2\) in the FULL ring,
Lagrange interpolation gives
\[
 C=\sum {W_i\over f'(U_i)},\qquad
 B+C(\beta+\chi)={W_\beta-W_\chi\over\beta-\chi}.
\tag{D17L}\label{eq:arithD17interpolation}
\]
The complete strip implies \(3<C<8\) and
\(|B+C(\beta+\chi)|<.174\).
At \(C=4,5,6,7\), only \(C=4,B=-7\) survives.
Then \(7\chi-4\chi^2\in(1.229,1.235)\),
so the small strip requires \(1.229<A<1.471\),
no integer.
For head eighteen the corresponding bounds are
\[
 \begin{gathered}
3.879<\alpha<3.880,\quad1.652<\beta<1.653,\quad
 .467<\chi<.469,\\
7.590<f'(\alpha)<7.605,\quad|f'(\beta)|>2.63,\quad
 f'(\chi)>4,\quad
2.120<\beta+\chi<2.121,\quad\beta-\chi>1.183 .
 \end{gathered}
\]
Interpolation gives \(2<C<6\),
\(|B+C(\beta+\chi)|<.199\).
For \(C=3,4,5\), the products lie in
\((6.360,6.363),(8.480,8.484),(10.600,10.605)\);
none is within \(.199\) of an integer.
No coefficient \(B\) exists.
Both cubic finite heads are excluded.

All FIVE finite heads have now been eliminated
at their ACTUAL fields.
Only NOW conclude that the head is DOTTED,
with \(4<U<11\), \(0<U_\sigma<2\).
Fork isolation and the dotted fork assumptions
proved conditionally in this subsection are NOT
exported to that remaining head.
We next prove them anew in the dotted case.

\subsection{Whole \(D_{17}\): degree reduction under a dotted head}
\label{app:arithD17degrees}

The complete fork table remains valid.
For an end helper, the integral \(v\) has norm
magnitude \(<6.607(.617)^{d-1}\).
Degree at least five gives norm \(<1\).
At degree four it is a unit;
parity~\eqref{eq:arithD17helperparity} then makes
\(Y=U/2\) integral.
With physical \(2<Y<11/2\) and three other
conjugates in \((0,1)\),
\[
                 0<|\mathrm N(Y(1-Y))|<99/256<1,
\]
impossible.
For an isolated fork the actual LONG+TWO
cycles above remain valid under a dotted head
BEFORE assigning its own type.
They again give integral \(W=V/2\), with
\(0<W<676/17\) and \(0<W_\sigma<4/17\)
at every other embedding.
The lower bound \(2<W\) is not used before
the finite-own exclusion below.
Degree at least four contradicts
\((676/17)(4/17)^3<1\).
Thus actual degree is two or three.
Rational degree was excluded by the named
WHOLE-form Meyer input.

We first remove cubic nonisolated forks.
The end-helper norm is \(\pm1,\pm2\).
Opposite signs of the two small roots give
magnitude \(<6.607(.617)(.382)<2\).
Two positive small roots in the far case,
or two negative in the near case, give magnitude
\(<1\).
Norm magnitude two can therefore occur only
as norm \(+2\) in the following two patterns.
For far roots \(t,-a,-b\), \(tab=2\),
\(0<a,b<.617\), the trace is five,
\(6<t<6.234\), and the integer second
coefficient is \(-6\) or \(-7\).
The polynomials
\(X^3-5X^2-6X-2\),
\(X^3-5X^2-7X-2\)
have discriminants \(-424,229\).
The first is not totally real.
At \(-.617\), the second has value
\(.180669887>0\), inconsistent with both
negative roots above that lower endpoint.
For near roots \(v,a,b\) with \(vab=2\),
the trace is seven, \(5.766<v<6\),
and the second coefficient is seven.
Its polynomial \(X^3-7X^2+7X-2\)
has discriminant \(-59\), impossible.
Thus a remaining cubic end-helper is a unit,
so \(U/2\) integral.

Here is the full EVEN-cubic-head conclusion,
also reused below.
For integral \(Y=U/2\), physical \(2<Y<11/2\)
and other roots in \((0,1)\) give
\[
                       0<|\mathrm N(Y(1-Y))|<99/64<2.
\]
Thus \(\mathrm N(Y)=1,\mathrm N(1-Y)=-1\).
Its polynomial is
\(X^3-TX^2+(T-1)X-1\), \(T=3,\ldots,7\).
The COMPLETE discriminants are
\(-23,-31,-23,49,257\).
The first three are not totally real;
at \(T=7\), evaluation at six is \(-1\),
forcing physical root \(>6\).
Only \(T=6\) remains.
Its reciprocal generates the cyclic field
\(K_7\), with full basis \(1,u,u^2\) for
\(u^3-5u^2+6u-1=0\).
The square discriminant gives cyclicity, so
choose its PHYSICAL largest root as generator.
The full-order proof and root/derivative bounds
are those just proved above and are algebraic
facts of this field, not a finite-head hypothesis.

For a cubic end-helper \(v=A+Bu+Cu^2\),
the full far/near strips and
\eqref{eq:arithD17interpolation} imply
\(-1<C<2\), hence \(C=0,1\), and
\(|B+C(\beta+\chi)|<.738\).
At \(C=0\), \(B=0\), and the small strip
forces \(A=0\), contrary to physical \(v\ge1\).
At \(C=1\), only \(B=-2\) is possible.
Since \(2\chi-\chi^2\in(.356,.359)\),
the small strip gives \(A=0\) in either role.
But \(v_\beta=\beta(\beta-2)<-.689\),
outside BOTH strips.
This exhausts every coefficient triple.
Cubic end helpers are impossible, so BOTH
actual cubic forks are isolated dotted.
Finite own in degree three would violate
\(\mathrm N(V)<4(8/17)^2<1\).
Each integral \(W\) satisfies
\eqref{eq:arithD17forkstrip},
with norm one or two; no equality of the two
fork parameters has been assumed.

\subsection{Whole \(D_{17}\): the entire isolated cubic field domain}
\label{app:arithD17cubic}

Choose ONE actual isolated cubic fork; the other
may have a different parameter.
Write its monic polynomial
\(P(X)=X^3-AX^2+BX-n\),
\(n=\mathrm N(W)=1,2\).
Its two other positive roots lie in
\((0,q)\), \(q=4/17\), physical
\(W<676/17\), hence integer trace \(A\le40\).
Since \(P(q)<0\),
\[
                   64-272A+1156B-4913n<0 .
\]
For \(n=2\), physical \(W>289/8\), and the
small-root arithmetic-geometric inequality gives
\(B\ge2\sqrt{2W}+2/W>17\).
The polynomial inequality gives
\(B<20642/1156<18\), impossible.
Thus \(W\) is a UNIT.

For the actual head/fork products in the
normalization \(r_H^2=2U,r_f^2=4W\),
the LONG and TWO cycles give integral \(g,L\),
with
\[
                   L-g=(8U-17)W-2U+4=:d_0,\qquad
                   g^2=2UW(g+d_0).
\]
At every dyadic prime, \(d_0\equiv-W\) is a unit.
The equation forces \(g\) into the prime and
then
\(2\operatorname{ord}(g)
 =\operatorname{ord}(2)+\operatorname{ord}(U)\).
If two is unramified this makes \(U/2\) integral.
The even-cubic-head conclusion above gives
actual field \(K_7\).
Apply the SAME full-order interpolation as for
finite head seven, now merely as algebra of the
field and the CURRENT unit-\(W\) strip.
Its complete \(C=4,\ldots,7\) tests leave only
\(C=4,B=-7\), then no integer \(A\).
Thus the unramified case is impossible, without
assuming the head itself is finite.

If two ramifies, the unit polynomial has a
repeated factor modulo two.
Its derivative \(X^2+B\) shows this requires
BOTH \(A,B\) odd.
The physical root and two small-root bounds give
\[
 \begin{gathered}
 A=19,21,\ldots,39,\qquad
 B^2>4(A-8/17),\qquad
 B<{272A+4849\over1156}.
 \end{gathered}
\]
The COMPLETE lower/upper odd-integer tests are
\[
\begin{array}{c|rrrrrrrrrrr}
 A&19&21&23&25&27&29&31&33&35&37&39\\ \hline
 B_{\min}&9&11&11&11&11&11&13&13&13&13&13\\
 B_{\max}&7&9&9&9&9&11&11&11&11&11&13 .
\end{array}
\]
Only
\[
                   P_{29}=X^3-29X^2+11X-1,\qquad
                   P_{39}=X^3-39X^2+13X-1
\]
remain.
Their discriminants are \(4(3)(383)\),
\(4(5021)\), with squarefree odd parts.
For an elementary full-order certificate,
trial primes below \(\sqrt{383}\) are
\(2,3,5,7,11,13,17,19\), none dividing383.
For5021, all primes below its square root are
\[
2,3,5,7,11,13,17,19,23,29,31,37,41,43,47,53,59,61,67;
\]
the respective remainders are
\[
1,2,1,2,5,3,6,5,7,4,30,26,19,33,39,39,6,19,63 .
\]
The shifted polynomials are two-Eisenstein:
\[
 P_{29}(X+1)=X^3-26X^2-44X-18,\qquad
 P_{39}(X+1)=X^3-36X^2-62X-26 .
\]
Only two can divide the power-order index.
At its uniformizer \(W-1\), valuations of
\(a_0+a_1\pi+a_2\pi^2\) are distinct modulo
three; a numerator not all divisible by two
has unique least valuation at most two.
Dividing by two cannot be integral.
Thus \(\mathcal O_K=\mathbb Z[W]\) in BOTH cases.

The exact small-root brackets for \(P_{29}\)
are \(.14<u<.15\), \(.23<v<.235\);
its physical root satisfies \(28.615<w<28.63\).
They follow from values
\(-.025656,.000875,.008067,-.003547125\)
at those four endpoints.
The physical derivative is \(>800\), both
small derivative magnitudes \(>2.27\).
For \(P_{39}\), values
\(-.040569,.000128,.019361,-.016952\)
give \(.11<u<.12\), \(.21<v<.22\),
\(38.66<w<38.68\), physical derivative \(>1480\),
small derivative magnitudes \(>3.45\).

Write the ACTUAL head \(U=a+bW+cW^2\) in
these proved FULL rings.
Its current strips are physical \((4,11)\),
other values \((0,2)\).
Lagrange interpolation bounds \(c\) respectively by
\[
             -2/2.27<c<11/800+2/2.27,\qquad
             -2/3.45<c<11/1480+2/3.45 .
\]
Both intervals are strictly inside \((-1,1)\),
so \(c=0\).
Then \(b=(U(w)-U(u))/(w-u)\) satisfies
\(0<b<11/(w-u)<1\), impossible for an integer.
Every ramified cubic case has been excluded.
The ACTUAL field is QUADRATIC.

\subsection{Whole \(D_{17}\): complete quadratic end-helper exclusion}
\label{app:arithD17quadratichelpers}

The dotted-head fork table already excludes
all four non-end white roles.
For the remaining far/near end-three roles
write \(v\)'s polynomial \(X^2-TX+n\).
The COMPLETE far ranges are
\(T=1,\ldots,6\), \(n=-3,-2,-1,1,2\);
near ranges \(T=2,\ldots,7\),
\(n=-2,-1,1,2,3,4\).
The quadratic formula and EXACT time signs leave
the following fourteen and fifteen rows:
\[
\begin{array}{c|c}
 (T,n,d)\text{ far}&(T,n,d)\text{ near}\\ \hline
(2,-1,2)&(3,-1,13)\\
(3,-1,13)&(3,1,5)\\
(3,-2,17)&(4,-1,5)\\
(4,-1,5)&(4,1,3)\\
(4,-2,6)&(4,2,2)\\
(4,1,3)&(5,-1,29)\\
(5,-1,29)&(5,-2,33)\\
(5,-2,33)&(5,1,21)\\
(5,-3,37)&(5,2,17)\\
(5,1,21)&(6,-1,10)\\
(6,-1,10)&(6,-2,11)\\
(6,-2,11)&(6,1,2)\\
(6,1,2)&(6,2,7)\\
(6,2,7)&(6,3,6)\\
 &(7,3,37).
\end{array}
\]
Here actual \(K=\mathbb Q(\sqrt d)\).
At each bounded pair \(T,n\), substitute the larger
root and its conjugate into the FULL strips.
Every missing pair is rational, nonreal or
violates an endpoint.
The safe far decimal overstrip admits one extra
\((3,1)\), \(v=\varphi^2\).
Its exact conjugate time numerator is
\((123-55\sqrt5)/2>0\), excluded by
\(123^2-5(55)^2=4\).
Thus the exact, rather than decimal, domain
is fully enforced.

For odd \(n\), \(v\) is a unit at EVERY dyadic
prime, so parity~\eqref{eq:arithD17helperparity}
makes \(U/2\) integral even with ramification.
The FULL quadratic head strip for \(Y=U/2\)
is \(2<Y<11/2\), \(0<Y'<1\).
Its integer trace is three through six and
norm \(1\le k\le T-2\); at trace six the
physical bound removes \(k=1,2\).
The COMPLETE eight actual heads and
\(h=-\mathrm N(U(U-2))\) are
\[
\begin{array}{c|c|c}
 d&U&h\\ \hline
5&2\varphi^2,\ 4\varphi^2&16,64\\
2&4+2\sqrt2&32\\
3&4+2\sqrt3&32\\
13&5+\sqrt{13}&48\\
17&5+\sqrt{17}&64\\
21&5+\sqrt{21}&48\\
6&6+2\sqrt6&96 .
\end{array}
\]
Odd-norm helper rows in fields29/10/37 are
immediately incompatible.
For all others, put
\(p=-16\mathrm N(H_c^2)>0\).
The complete time-norm tests are
\[
\begin{array}{c|c|r|l}
 \text{role}&(T,n,d)&p&h\\ \hline
 \text{far}&(2,-1,2)&287&32\\
 &(3,-1,13)&711&48\\
 &(4,-1,5)&1271&16,64\\
 &(4,1,3)&423&32\\
 &(5,1,21)&983&48\\
 &(6,1,2)&1679&32\\
 \text{near}&(3,-1,13)&207&48\\
 &(3,1,5)&311&16,64\\
 &(4,-1,5)&599&16,64\\
 &(4,1,3)&839&32\\
 &(5,1,21)&1503&48\\
 &(6,1,2)&2303&32\\
 &(6,3,6)&503&96 .
\end{array}
\]
All \(p/(16h)\) are nonsquares.
At \(h=32,96\) the odd numerator gives odd
two-adic denominator valuation.
The four \(h=48\) reduced cases are
\(237/256,983/768,69/256,501/256\);
the middle denominator has odd three valuation,
the other numerators are nonsquares.
For \(h=16,64\), the numerators
\(1271,311,599\) are nonsquares.
These are NECESSARY actual common-time-line
square tests, not sufficiency of a norm-square.

For even-norm rows DO NOT assume \(U/2\)
integral.
Their complete actual fields, full head strips
and \(h\)'s are
\[
\begin{array}{c|l|l}
 d&U&h\\ \hline
2&3+\sqrt2,\ 3+2\sqrt2,\ 4+2\sqrt2,\
     5+3\sqrt2,\ 6+3\sqrt2&7,7,32,63,36\\
6&3+\sqrt6,\ 4+\sqrt6,\ 5+2\sqrt6,\ 6+2\sqrt6&
     15,20,15,96\\
7&3+\sqrt7,\ 4+\sqrt7&12,27\\
11&4+\sqrt{11},\ 5+\sqrt{11}&35,28\\
17&2+\eta,\ 3+\eta,\ 4+2\eta,\ 5+2\eta&
      8,16,64,19\\
33&3+\xi,\ 4+\xi&24,24 .
\end{array}
\]
Here \(\eta=(1+\sqrt{17})/2\),
\(\xi=(1+\sqrt{33})/2\).
The FULL integral bases and \(0<U'<2,4<U<11\)
give coefficient bounds respectively
\(3,2,2,1,2,1\).
At fixed positive coefficient, the width-two
conjugate interval gives two consecutive constants;
the physical strip removes the excess.
For example at \(\sqrt7\), coefficient two
already has least physical \(6+2\sqrt7>11\).
Thus no full-order head was omitted.

The COMPLETE remaining even-norm time rows are
\[
\begin{array}{c|c|r|l}
 \text{role}&(T,n,d)&p&\text{all }h\\ \hline
 \text{far}&(3,-2,17)&200&8,16,64,19\\
 &(4,-2,6)&828&15,20,15,96\\
 &(5,-2,33)&1592&24\\
 &(6,-2,11)&2492&35,28\\
 &(6,2,7)&252&12,27\\
 \text{near}&(4,2,2)&92&7,7,32,63,36\\
 &(5,-2,33)&72&24\\
 &(5,2,17)&824&8,16,64,19\\
 &(6,-2,11)&668&35,28\\
 &(6,2,7)&1692&12,27 .
\end{array}
\]
Every entry follows by expanding
\[
16\mathrm N(H_c^2)
 =(17T\pm4)^2n+(17T\pm4)(-17n-4)T+(-17n-4)^2 .
\]
Every ratio \(p/(16h)\) is nonsquare EXCEPT
far \((T,n,d)=(3,-2,17)\), \(h=8\).
The far nonexceptional rows retain odd valuations
at \(23,199,89,7\).
In the exceptional field the other ratios are
\(25/32,25/128,25/38\), nonsquares.
The near rows retain odd valuations at
\(23,3,103,167,47\).
Thus a genuine square-norm exception is kept.

For that exception
\[
 t=(3+\sqrt{17})/2,\quad U=(5+\sqrt{17})/2,
\]
the ACTUAL time squares are
\[
 H_c^2=(225+55\sqrt{17})/8,\qquad
 H_0^2=(11+3\sqrt{17})/2 .
\]
Their ratio
\(R=(165+35\sqrt{17})/64\) has norm \(25/16\).
If \(R=z^2\) in \(K\), then
\(\mathrm N(z)=\pm5/4\), and
\[
                    \operatorname{Tr}(z)^2=245/32
                                      \quad\hbox{or}\quad85/32 .
\]
Neither is a rational square, by its odd
two-adic denominator valuation.
But the trace is rational.
The last actual common-time square fails.
All quadratic end-helper rows are now excluded,
in both signed orientations.

Both ACTUAL forks are therefore isolated under
the dotted quadratic head.
Their original LONG+TWO cycles, applied BEFORE
assuming own dotted, give integral \(W=V/2\)
with \(0<W<676/17\) and \(0<W'<4/17\).
The lower bound \(2<W\) is deferred until
finite owns have been excluded.
Finite own would have \(0<W<2\), other
\(0<W'<4/17\), hence nonzero integer norm
\(0<\mathrm N(W)<8/17\), impossible.
Parallel has its actual null pair.
Both forks are isolated DOTTED, without
equating their two \(W\) parameters.

\subsection{Whole \(D_{17}\)'s distinct actual inclusion and final budget}
\label{app:arithD17domain}

Choose ONE of the two actual isolated dotted forks.
Its raw square is \(4W\), sole old product
\(-2W\), spatial \(W(P_{16}\pm e_{17})\),
time square \(W(17W-4)\).
The actual head has square \(2U\), spatial
\(Ue_1\), time square \(U(U-2)\).
The original \(K\)-time line gives
\[
 R_0=U(U-2)W(17W-4)\in K^2,\qquad
 g=UW-\sqrt{R_0},
\]
with PHYSICAL positive square root.
The original LONG path and TWO cycles give
\[
                   g,\ L={g^2\over2UW}\in\mathcal O_K .
\]
All physical/conjugate strips are now
\[
 K=\mathbb Q(U)=\mathbb Q(W),\quad
 4<U<11,\quad0<U'<2,\quad
 2<W<676/17,\quad0<W'<4/17 .
\]
The actual independent conjugate root lines give
\(0<L'<4\) when \(g\ne0\).
Here the current physical strips force
\(d_0=(8U-17)W-2U+4>0\), hence \(g<0\),
so no orthogonal exception is silently lost.
The true physical mutual \(L\) is finite
native or dotted, while parallel four has
its forbidden actual null pair.
Thus EVERY predicate of
Proposition~\ref{prop:pair-output}(iii) holds
for this ORIGINAL D17 geometry.

This inclusion is separate from
Section~\ref{app:arithB17domain}:
D17 raw fork norm is \(4W\);
B17 raw short norm was \(V\).
No equation of different root scales or
arithmeticity of a face supplied either inclusion.
The existing full-order algorithm covers every
\(\Delta<121\), every strict-strip model and
ordered pair with both field-square norm signs.
Its 36-field/120-pair output, explicitly ASSUMED
as the recorded computational proposition, has
the same TWO necessary survivors.
The field/counter/survivor combine does not
independently recompute the118 rejected tests.

For either actual survivor, use its actual
head point and ONE chosen signed fork.
The head gain at cost \(.072\) from \(P_5\)
is \(>.043\), by the full rational seals already
printed in Section~\ref{app:arithB17domain}.
The actual fork has
\(C_f^2=17W/4>442/5\), giving
\[
                 x\cdot(P_{16}\pm e_{17})>3.7 .
\]
Pay \(.13\) from the DISTINCT \(P_{14}\) price:
\[
             y=w_{14}+.36(Z-P_5/5)
                          +.13(P_{16}\pm e_{17}-P_{14}).
\]
The signed residual has norm three and
vanishes on all reference coordinates;
it is orthogonal to the head residual.
Both individual payments are legal BEFORE
using their original floors.
ONE square on the SAME original \(x\) gives
\[
              \|x\|^2>.9084+.043+.26(3.7-3.311)-.0169(3)
                                     =1.00184>1 .
\]
Every actual field and actual fork type has
now been covered.
This closes whole \(D_{17}\) at the stated
external and explicitly computed premises.
The argument needed one chosen fork for the
last contradiction, while preserving both
signed walls and all original constraints.



